\section{A104722: a recurrence for repeated Catalan numbers}

Write $C_r=\frac1{r+1}\binom{2r}{r}$ for the Catalan numbers, and
define the repeated sequence and its self-convolution by
\begin{equation}\label{eq:a104722-source}
 b_n=C_{\lfloor n/2\rfloor},\qquad
 a_n=\sum_{k=0}^{n}b_kb_{n-k}\qquad(n\ge0).
\end{equation}
This is the literal definition of A104722 \cite{a104722-source}.
The initial values are $a_0,a_1,a_2,a_3=1,2,3,4$.

\begin{theorem}\label{thm:a104722}
For every integer $n\ge4$,
\[
 (n+4)a_n+(n+1)a_{n-1}
 -4(n+1)a_{n-2}+4(2-n)a_{n-3}=0.
\]
At $n=3$, the left side is $4$, so the starting index cannot be
lowered to $3$.
\end{theorem}

\begin{proof}
The standard Catalan convolution and ratio identities are
\begin{equation}\label{eq:a104722-catalan}
 \sum_{j=0}^{r}C_jC_{r-j}=C_{r+1},\qquad
 (r+2)C_{r+1}=2(2r+1)C_r.
\end{equation}
The first is the coefficient form of $C(x)=1+xC(x)^2$;
the second follows from the binomial formula.

For the odd argument $2r+1$, pair the even and odd summation
indices in \eqref{eq:a104722-source}. Both terms in the $j$th
pair are $C_jC_{r-j}$, giving
\begin{equation}\label{eq:a104722-odd}
 a_{2r+1}=2\sum_{j=0}^{r}C_jC_{r-j}=2C_{r+1}.
\end{equation}
For the even argument $2r+2$, the even terms sum to
$\sum_{j=0}^{r+1}C_jC_{r+1-j}=C_{r+2}$ and the odd terms sum to
$\sum_{j=0}^{r}C_jC_{r-j}=C_{r+1}$. Hence
\begin{equation}\label{eq:a104722-even}
 a_{2r+2}=C_{r+2}+C_{r+1}.
\end{equation}
Both identities hold for every $r\ge0$.

Let $R(n)$ denote the recurrence's left side and set
\[
 H_r=(r+2)C_{r+1}-2(2r+1)C_r.
\]
The second identity in \eqref{eq:a104722-catalan} says $H_r=0$.
Substitute \eqref{eq:a104722-odd} and \eqref{eq:a104722-even}
into $R(n)$ and collect the three adjacent Catalan numbers.
For the two parity cases this gives the exact certificates
\begin{align*}
 R(2r+4)&=2H_{r+2}+6H_{r+1},\\
 R(2r+5)&=6H_{r+2}+2H_{r+1}.
\end{align*}
For clarity, their expanded forms are
\[
\begin{aligned}
R(2r+4)
 &=(2r+8)C_{r+3}+(-2r-2)C_{r+2}-(24r+36)C_{r+1},\\
R(2r+5)
 &=(6r+24)C_{r+3}-(22r+54)C_{r+2}-(8r+12)C_{r+1}.
\end{aligned}
\]
Thus both parity cases vanish and exhaust all integers $n\ge4$.
At the preceding index, direct substitution gives
\[
 R(3)=7\cdot4+4\cdot3-16\cdot2-4\cdot1=4.
\]
\end{proof}

\paragraph{Prior work and scope.}
The source records Mathar's conjecture on November 9, 2012, but
prints no lower bound on $n$ \cite{a104722-source}; the theorem
supplies the necessary bound. This is a single recurrence family.
The Catalan expressions are established background. The OEIS
entry credits James Sellers's May 19, 2008 generating-function
argument for agreement with A059348 after the first term.
The related sequence has a different OEIS offset
\cite{a104722-related}. Zhao, Li and Sun's Theorem 3.7
\cite{a104722-zhao} gives a related permutation enumeration in
Catalan numbers; no novelty is claimed here for those formulas.

The contribution of this package is the explicit recurrence
derivation from the literal convolution, its endpoint correction,
and its checked Lean proof. The bounded prior-work search is
documented separately and does not establish global priority.
The formalization uses mathlib's natural-number Catalan sequence,
proves both parity formulas by finite-sum pairing, derives the ratio
from the central-binomial-coefficient identity, and checks both
linear-combination certificates. Its final results are
\texttt{A104722.source\_recurrence} and
\texttt{A104722.index\_three\_exception}.
